\documentclass[11pt]{article}
\usepackage[T1]{fontenc}
\usepackage[utf8]{inputenc}
\usepackage{lmodern}
\usepackage{amsmath,amssymb,amsthm,mathtools}
\usepackage{booktabs,array,microtype}
\usepackage[a4paper,margin=27mm]{geometry}
\usepackage[hidelinks]{hyperref}
\usepackage{xcolor}

\newtheorem{theorem}{Theorem}[section]
\newtheorem{lemma}[theorem]{Lemma}
\newtheorem{proposition}[theorem]{Proposition}
\newtheorem{corollary}[theorem]{Corollary}
\theoremstyle{definition}
\newtheorem{definition}[theorem]{Definition}
\theoremstyle{remark}
\newtheorem{remark}[theorem]{Remark}

\newcommand{\R}{\mathbb{R}}
\newcommand{\Z}{\mathbb{Z}}
\newcommand{\cA}{\mathcal{A}}
\newcommand{\cC}{\mathcal{C}}
\newcommand{\Cov}{\operatorname{Cov}}
\newcommand{\TV}{\operatorname{TV}}
\newcommand{\diag}{\operatorname{diag}}
\newcommand{\lean}[1]{\nolinkurl{#1}}
\newcommand{\ModuleHash}{\nolinkurl{a3089a961a3f7659c9a5317981168ffcccd980bffc10f894adab1788076ed43b}}
\newcommand{\OleanHash}{\nolinkurl{739a3c64a53802e1f6483ce63b53f72b1ce06cdbf0ddd856b6efbb70f7f5ffa7}}
\newcommand{\BridgeHash}{\nolinkurl{1dd660dceb92b133a5dd75e4e9a7bb3afad9d5d635010c61b8586c686227cbe9}}
\newcommand{\BridgeOleanHash}{\nolinkurl{77eb1ed915efc4c3ce7905de7b72bc3d82e117d72d18a8bab21df2da35da9555}}
\setlength{\emergencystretch}{2em}

\title{Exact Diagonal and Mixed Reconstructed Correlator Decay\\
from a Dobrushin Window for Finite Coupled $\Z_2$ Slices}
\author{Lluis Eriksson}
\date{Technical Note\\4 August 2026}

\begin{document}
\maketitle

\begin{abstract}
We give one end-to-end account of a finite-volume chain that had previously
been split between a Dobrushin comparison development and a reconstruction
interface.  For anisotropic coupled $\Z_2$ slices, the explicit condition
\[
  0<\alpha<1,
  \qquad 2\tanh|\beta|+2\tanh|\gamma|\leq\alpha
\]
first controls the intrinsic single-site interdependence matrix, then the
finite-strip covariances, then the fluctuation norm of the normalized symmetric
transfer kernel.  An exact square-root dressing transports that bound to the
reconstructed site form.  Consequently, with $m=-\log\alpha>0$, one obtains
for every spatial extent and every finite temporal separation
\[
  |C^{\rm rec}_L(u;n)|\leq \|u\|_2^2 e^{-mn}.
\]
The exponent and prefactor are unchanged by reconstruction.  This version
also defines the mixed correlator in both abstract and reconstructed
coordinates, proves their exact equality, and machine-checks
\[
  |B^{\rm rec}_L(u,v;n)|\leq \|u\|_2\|v\|_2 e^{-mn}
\]
including the separate zero-time case.  Here $u$ and $v$ range over the full
real algebra of functions on a boundary slice.  No sharper Dobrushin
window, thermodynamic limit, continuum reconstruction, or Yang--Mills mass gap
is claimed.  The analytic estimate is inherited; the new formal increment is
the mixed observable endpoint and its inclusion in the uniform theorem.
\end{abstract}

\section{Result and exact scope}

Let
\[
  X_L=(\operatorname{Fin}(L+1)\to\operatorname{Fin}2),
  \qquad \cA_L=\R^{X_L}.
\]
The space $\cA_L$ is not merely a chosen vector class: with pointwise
multiplication it is the complete finite-dimensional commutative algebra of
real observables on one boundary slice.  Thus every cylinder observable that
depends on that whole slice, as well as every finite sum and pointwise product
of such observables, is represented by some $u\in\cA_L$.  What remains finite
is the slice itself and the temporal separation.

The slice weight $w_{\gamma,L}:X_L\to\R_{>0}$ contains the spatial coupling
$\gamma$.  The temporal kernel $K_\beta$ is dressed symmetrically as
\[
  S_{w,\beta}=\diag(\sqrt w)K_\beta\diag(\sqrt w).
\]
For positive Perron data $(\lambda_L,\Omega_L)$ put
\[
  P_L=\lambda_L^{-1}S_{w_{\gamma,L},\beta},
  \qquad P_L\Omega_L=\Omega_L,
  \qquad \|\Omega_L\|_2=1.
\]
The reconstructed site form and forced transfer operator are denoted by
$\mathfrak q_w$ and $T_{w,\beta}$.  Their precise finite sums appear in
Section~\ref{sec:reconstruction}.

\begin{theorem}[end-to-end finite reconstructed clustering]
\label{thm:main}
Assume
\[
  0<\alpha<1,
  \qquad 2\tanh|\beta|+2\tanh|\gamma|\leq\alpha .
\]
Set $m=-\log\alpha$.  Then $m>0$ and, for every $L\geq0$, there are positive
normalized Perron data $(\lambda_L,\Omega_L)$ such that
\[
  \|P_L-|\Omega_L\rangle\langle\Omega_L|\|\leq e^{-m}
\]
and, for all $u\in\cA_L$ and $n\in\mathbb N$,
\[
  |C^{\rm rec}_L(u;n)|
  \leq \|u\|_2^2 e^{-mn}.
\]
For all $u,v\in\cA_L$ and $n\in\mathbb N$ one also has
\[
  |B^{\rm rec}_L(u,v;n)|
  \leq \|u\|_2\|v\|_2 e^{-mn}.
\]
Moreover $C^{\rm rec}_L(u;n)$ is exactly the Euclidean connected correlator
of $P_L$, and $B^{\rm rec}_L(u,v;n)$ is exactly its mixed connected matrix
element, not merely bounded by it.
\end{theorem}

The binder order is load-bearing:
\[
  \exists m>0\;\forall L\;\exists(\lambda_L,\Omega_L)\;
  \forall u,v,n.
\]
Perron data may depend on the volume; the exponent does not.  The theorem is
the mathematical statement exposed by
\lean{os\_reconstruction\_uniform\_gap}.  The following sections give the
entire implication chain rather than referring to the analytic input as a
black box.

\section{The intrinsic Dobrushin estimate}
\label{sec:dobrushin}

Let $I$ be a finite site set, $S=\{\pm1\}$, and let $W:S^I\to\R_{>0}$ be a
strictly positive Gibbs weight.  Its normalized measure and heat-bath
conditional are
\[
  \mu_W(\eta)=\frac{W(\eta)}{\sum_\xi W(\xi)},
  \qquad
  p_i^W(\eta,s)=
  \frac{W(\eta^{i\to s})}{\sum_{t\in S}W(\eta^{i\to t})}.
\]
Define the intrinsic interdependence coefficient
\[
  c^W_{ik}=
  \max_{\eta\equiv\eta'\;\mathrm{off}\;k}
  \|p_i^W(\eta,\cdot)-p_i^W(\eta',\cdot)\|_{\TV}.
\]
For an observable $f$ let
\[
  \delta_i(f)=\max_{\eta\equiv\eta'\;\mathrm{off}\;i}
  |f(\eta)-f(\eta')|.
\]

\begin{lemma}[finite Dobrushin comparison]
\label{lem:comparison}
If $\sum_k c^W_{ik}\leq\alpha<1$ for every $i$, then
\[
  |\Cov_{\mu_W}(f,g)|
  \leq \frac14\sum_{i,j}\delta_i(f)
       \left(\sum_{q\geq0}(c^W)^q_{ij}\right)\delta_j(g).
\]
If $c^W_{ik}=0$ whenever $d(i,k)>1$, then
\[
  |\Cov_{\mu_W}(f,g)|
  \leq \frac{1}{4(1-\alpha)}
  \sum_{i,j}\delta_i(f)\delta_j(g)\alpha^{d(i,j)}.
\]
\end{lemma}

\begin{proof}
The heat-bath update $E_i$ is invariant under $\mu_W$.  Indeed, after
expanding $\mathbb E_{\mu_W}[E_iF]$ over $(\eta,s)$, the involution
$(\eta,s)\mapsto(\eta^{i\to s},\eta_i)$ preserves the local partition sum
and cancels its denominator.  The conditional Gr\"uss estimate contributes
the factor $1/4$.  A single update transports the oscillation vector by
$c^W$; iteration and invariance give the Neumann resolvent
$D=\sum_{q\geq0}(c^W)^q$.  The row-sum hypothesis makes the series
entrywise summable.  Under range one, $(c^W)^q_{ij}=0$ for
$q<d(i,j)$ and its row sum is at most $\alpha^q$, hence
\[
  D_{ij}\leq\sum_{q\geq d(i,j)}\alpha^q
  =\frac{\alpha^{d(i,j)}}{1-\alpha}.
\]
Substitution proves both bounds.  The corresponding checked endpoints are
\lean{gibbs\_covar\_le\_resolvent} and
\lean{gibbs\_covar\_exp\_decay}.
\end{proof}

\section{The anisotropic rectangle and the common rate}
\label{sec:rectangle}

For a symmetric zero-diagonal Ising coupling $J$, write
\[
  W_J(\eta)=\exp\!\left(\frac12\sum_{a,b}J_{ab}\eta_a\eta_b\right).
\]
The heat-bath conditional is the sigmoid of the local field
$h_i(\eta)=\sum_kJ_{ik}\eta_k$:
\[
  p_i^{W_J}(\eta,+1)=\frac{1+\tanh h_i(\eta)}2.
\]
Changing spin $k$ moves $h_i$ by at most $2|J_{ik}|$.  The elementary
two-point sigmoid calculation therefore gives
\[
  c^{W_J}_{ik}\leq\tanh|J_{ik}|.
\]
This is an envelope, not a claim that the intrinsic coefficient always
attains it.

Consider a finite rectangular strip, with coupling $\beta$ in the temporal
direction and $\gamma$ in the spatial direction, and free boundary.  Every
site has at most two neighbors of each type, hence
\[
  \sum_k c^{W_J}_{ik}
  \leq 2\tanh|\beta|+2\tanh|\gamma|
  \leq\alpha.
\]
Boundary sites have a smaller row sum.  Since the coupling has range one for
Manhattan distance, Lemma~\ref{lem:comparison} yields, uniformly in both side
lengths,
\[
  |\Cov_{\mu_{L,T}}(f,g)|
  \leq \frac{1}{4(1-\alpha)}
  \sum_{p,q}\delta_p(f)\delta_q(g)\alpha^{d_\square(p,q)}.
\]
For observables supported on the two end slices of a strip of temporal length
$n$, every contributing pair has $d_\square(p,q)\geq n$.  Therefore
\begin{equation}
\label{eq:end-slice-decay}
  |\Cov_{\mu_{L,n}}(f,g)|
  \leq C_{L,f,g}\,\alpha^n,
\end{equation}
where $C_{L,f,g}$ may depend on the spatial extent and the observables, but not
on $n$.  This precise freedom is sufficient for the operator transport.  The
checked rectangle endpoints are \lean{rect\_ising\_uniform\_decay} and
\lean{rect\_ising\_uniform\_two\_point}.

\section{From strip covariances to an operator gap}
\label{sec:transport}

Let $M$ be a symmetric matrix on a finite set $X$, with a strictly positive
vector $\Omega$ satisfying $M\Omega=\Omega$ and $\|\Omega\|_2=1$.  For a path
$v=(v_0,\ldots,v_n)$ define
\[
  W_n(v)=\Omega(v_0)\prod_{k<n}M_{v_kv_{k+1}}\Omega(v_n).
\]
Successively summing the eigen-relation shows that $\sum_vW_n(v)=1$.  For
endpoint observables $f,g:X\to\R$, direct matrix multiplication gives
\begin{equation}
\label{eq:band-identity}
  \operatorname{bandCov}_n(f,g)
  =\langle(f-\mathbb Ef)\Omega,
      M^n(g-\mathbb Eg)\Omega\rangle.
\end{equation}
No spectral positivity is used; symmetry and the Perron relation suffice.

For the physical slice weight $w$ and temporal kernel $K_\beta$, take
\[
  S=\diag(\sqrt w)K_\beta\diag(\sqrt w),
  \qquad M=S/\lambda,
  \qquad \psi=\Omega/\sqrt w.
\]
Then the band weight is exactly a free-strip Gibbs weight with the two endpoint
factors $\psi$, divided by $\lambda^n$.  The identity is pointwise, so the
boundary factors change only $C_{L,f,g}$ in
\eqref{eq:end-slice-decay}; they do not change the common rate $\alpha$.
The currying bijection between rectangle configurations and slice paths then
identifies the free strip covariance with the rectangle covariance.

\begin{proposition}[uniform projected norm]
\label{prop:gap}
Under the window of Theorem~\ref{thm:main}, for every $L$ the normalized
symmetric slice kernel has Perron data with
\[
  \|M_L-|\Omega_L\rangle\langle\Omega_L|\|\leq\alpha.
\]
\end{proposition}

\begin{proof}
Equation~\eqref{eq:end-slice-decay}, the exact tilt identity, and
\eqref{eq:band-identity} give, for every vector $v$, a constant $C_{L,v}$ with
\[
  |\langle v,(M_L^n-|\Omega_L\rangle\langle\Omega_L|)v\rangle|
  \leq C_{L,v}\alpha^n.
\]
The fluctuation operator is self-adjoint.  If its norm exceeded $\alpha$, a
finite-dimensional spectral vector in the fluctuation sector would make the
left side decay at a strictly slower exponential rate, contradicting the
display.  Hence its norm is at most $\alpha=e^{-m}$ with
$m=-\log\alpha$.  In Lean this is the composition ending at
\lean{dobrushin\_ising\_uniform\_gap}, through
\lean{abstract\_uniform\_gap}.
\end{proof}

This step explains why a volume-dependent prefactor in the measure estimate
does not spoil the conclusion: the contradiction is performed separately in
each finite volume, whereas the forbidden rate is common.

\section{Exact reconstruction}
\label{sec:reconstruction}

For $u,v\in\cA_L$ define the reconstructed site form
\[
  \mathfrak q_w(Q_wu,Q_wv)=\sum_{\sigma\in X_L}u(\sigma)v(\sigma),
  \qquad (Q_wu)(\sigma)=\sqrt{w(\sigma)}u(\sigma).
\]
In the implementation the form is first written with the weight $w$ and the
inverse square-root dressing; strict positivity of $w$ cancels those factors
exactly.  Consequently $Q_w$ is an isometry from the Euclidean boundary space
onto the real reconstructed site coordinates.

The forced transfer matrix satisfies the exact intertwining relation
\begin{equation}
\label{eq:intertwining}
  T_{w,\beta}Q_wu=Q_wS_{w,\beta}u.
\end{equation}
Induction gives the equality for every finite iterate.  After Perron
normalization,
\[
  (\lambda^{-1}T_{w,\beta})^nQ_wu=Q_wP^nu.
\]
The checked declarations are \lean{siteForm\_qEmbed},
\lean{transferOp\_qEmbed}, \lean{transferOp\_iterate\_qEmbed}, and
\lean{transferOp\_qEmbed\_tilt\_iterate}.

Define the reconstructed connected autocorrelator by
\[
  C^{\rm rec}_{w,\beta,\lambda,\Omega}(u;n)
  =\operatorname{Re}\mathfrak q_w
     \left(Q_wu,(\lambda^{-1}T_{w,\beta})^nQ_wu\right)
   -\operatorname{Re}\mathfrak q_w(Q_w\Omega,Q_wu)^2.
\]
Equations above give the exact identity
\begin{equation}
\label{eq:corr-exact}
  C^{\rm rec}_{w,\beta,\lambda,\Omega}(u;n)
  =\langle u,P^nu\rangle-\langle\Omega,u\rangle^2.
\end{equation}
Thus, writing $\Pi=|\Omega\rangle\langle\Omega|$ and
$R_n=P^n-\Pi$,
\[
  |C^{\rm rec}(u;n)|=|\langle u,R_nu\rangle|
  \leq\|u\|_2^2\|P-\Pi\|^n.
\]
For $n\geq1$ this uses $P^n-\Pi=(P-\Pi)^n$; for $n=0$ it uses
$\|I-\Pi\|\leq1$.  This separate zero-time case is necessary because the
zeroth power of $P-\Pi$ is $I$, not $I-\Pi$.
Combining this with Proposition~\ref{prop:gap} proves
Theorem~\ref{thm:main}.  The endpoint is checked as
\lean{reconstructedConnCorr\_eq\_connCorr} and
\lean{reconstructedConnCorr\_decay}.

\section{The full boundary observable algebra}
\label{sec:observables}

The diagonal notation in \eqref{eq:corr-exact} can obscure the scope.  For
$u,v\in\cA_L$ define the mixed connected form
\[
  B^{\rm rec}_L(u,v;n)
  =\langle u,P_L^nv\rangle
   -\langle\Omega_L,u\rangle\langle\Omega_L,v\rangle.
\]
It is bilinear and, since $P_L$ is symmetric, symmetric in $u,v$.  The
implementation now contains the abstract definition \lean{mixedConnCorr} and
the reconstructed definition \lean{reconstructedMixedConnCorr}; the latter is
written directly with the two reconstructed boundary vectors rather than
obtained by a paper-only polarization step.

\begin{corollary}[mixed boundary observables]
\label{cor:mixed}
Under the hypotheses of Theorem~\ref{thm:main},
\[
  |B^{\rm rec}_L(u,v;n)|
  \leq\|u\|_2\|v\|_2e^{-mn}
\]
for all $u,v\in\cA_L$ and $n\in\mathbb N$.
\end{corollary}

\begin{proof}
Put $\Pi_L=|\Omega_L\rangle\langle\Omega_L|$.  For $n\geq1$, the vacuum
projection commutes with $P_L$ and
$P_L^n-\Pi_L=(P_L-\Pi_L)^n$.  Cauchy--Schwarz and
$\|P_L-\Pi_L\|\leq e^{-m}$ give the claim.  At $n=0$, use
$\|I-\Pi_L\|\leq1=e^{-m\cdot0}$.
\end{proof}

The positive-time identity is checked by
\lean{VacuumTransfer.mixedConnCorr\_eq}.  The zero-time contraction
$\|I-\Pi_L\|\leq1$ is now a named theorem,
\lean{norm\_one\_sub\_vacuumProjection\_le}; both cases are assembled by
\lean{mixed\_clustering\_of\_gap}.  Exact transport to reconstructed
coordinates is \lean{reconstructedMixedConnCorr\_eq\_mixedConnCorr}, and the
product-norm estimate is \lean{reconstructedMixedConnCorr\_decay}.  Finally,
\lean{os\_reconstruction\_uniform\_gap} exports the mixed clause under the
same volume-independent exponent as the diagonal clause.

This formal endpoint shows that the result is not restricted to one preferred
scalar observable.  Every real boundary function is admitted, and any finite
covariance matrix built from a family $u_1,\ldots,u_k\in\cA_L$ inherits the
same entrywise exponential rate.  It does not enlarge the support beyond one
finite slice or construct an infinite-time observable algebra.

\section{Machine-checked map and reproducibility}

This version changes two Lean modules.  Relative to v3 it adds two definitions
and five named theorems: the abstract and reconstructed mixed correlators, the
positive-time structural identity, the zero-time projection contraction, the
mixed gap-to-clustering estimate, exact reconstruction, and reconstructed
mixed decay.  It also strengthens the uniform endpoint by adding the mixed
clause without changing its Dobrushin hypotheses.

The exact reconstruction source has $11621$ LF bytes and SHA-256
\begin{center}\ModuleHash\end{center}
and its materialized \texttt{.olean} has SHA-256
\begin{center}\OleanHash\end{center}
The abstract bridge source has $36484$ LF bytes and SHA-256
\begin{center}\BridgeHash\end{center}
and its materialized \texttt{.olean} has SHA-256
\begin{center}\BridgeOleanHash\end{center}

Validation used Colab Pro+ with CPU and high RAM; no local Lean or Lake command
was run.  The target build completed $8191$ jobs with exit code zero.  The
dependency build completed $8482$ jobs with exit code zero.  Four exhaustive
oracle chunks were derived from the exact local \texttt{oracle\_check.lean}
body and emitted $3054$ reports: $3027$ used only Lean's standard axioms
\lean{Classical.choice}, \lean{Quot.sound}, and \lean{propext}; $27$ were
axiom-free; there were zero nonstandard axioms, zero \lean{sorryAx}, and zero
core errors.  The runtime was then disconnected and deleted.  A frozen-object
manifest records the body commit, LF/CRLF projections, build logs, oracle
hashes, and PDF hashes for independent assessment.

\begin{center}\small
\begin{tabular}{>{\raggedright\arraybackslash}p{0.40\linewidth}
                >{\raggedright\arraybackslash}p{0.50\linewidth}}
\toprule
Mathematical step & Checked endpoint or source lane\\
\midrule
Intrinsic heat-bath comparison & \lean{gibbs\_covar\_le\_resolvent}\\
Range-one exponential decay & \lean{gibbs\_covar\_exp\_decay}\\
Anisotropic rectangle, all volumes & \lean{rect\_ising\_uniform\_decay}\\
Band covariance identity & \lean{band\_covariance\_eq}\\
Measure decay to operator norm & \lean{abstract\_uniform\_gap}\\
Perron tilt and rectangle feed & \lean{dobrushin\_ising\_uniform\_gap}\\
Exact reconstructed correlator & \lean{reconstructedConnCorr\_eq\_connCorr}\\
Reconstructed decay & \lean{reconstructedConnCorr\_decay}\\
Exact mixed reconstructed correlator & \lean{reconstructedMixedConnCorr\_eq\_mixedConnCorr}\\
Mixed reconstructed decay & \lean{reconstructedMixedConnCorr\_decay}\\
Uniform physical endpoint & \lean{os\_reconstruction\_uniform\_gap}\\
\bottomrule
\end{tabular}
\end{center}

The analytic declarations in the first six rows predate the reconstruction
module.  This paper does not relabel them as new; it reproduces their argument
so the endpoint can be read without consulting a separate paper.

\section{Limitations and precise contribution}

The proof is finite-volume throughout.  It does not construct an
infinite-volume Gibbs state, prove boundary-condition independence, take a
thermodynamic or continuum limit, reconstruct Wightman fields, or prove a
Yang--Mills mass gap.  The Dobrushin window is sufficient and deliberately
conservative; for the two-dimensional Ising cell it is weaker than the exact
solution.  No improvement of that window or of $m=-\log\alpha$ is claimed.

What the integrated statement adds is exact compositional control.  The
classical single-site condition is carried through a finite rectangle, a
Perron boundary tilt, a projected operator norm, and the reconstructed site
form.  Every transition is an equality or a bound with its quantifiers visible,
and reconstruction loses neither rate nor prefactor.  The mixed-observable
theorem makes explicit---in the formal interface, not only in prose---that two
independent boundary variables range over the full real function algebra of
the finite slice.  This is a genuine formal increment but no new analytic
estimate: it packages the inherited operator-norm gap through exact
reconstruction with the sharp product prefactor.  These are useful endpoints
of a formalization program, but they remain finite and high-temperature.

\section*{Primary references}

\begin{thebibliography}{9}

\bibitem{Dobrushin1968}
R.~L. Dobrushin,
\emph{The description of a random field by means of conditional probabilities
and conditions of its regularity},
Theory of Probability and its Applications \textbf{13} (1968), 197--224.
\href{https://doi.org/10.1137/1113026}{doi:10.1137/1113026}.

\bibitem{Dobrushin1970}
R.~L. Dobrushin,
\emph{Prescribing a system of random variables by conditional distributions},
Theory of Probability and its Applications \textbf{15} (1970), 458--486.
\href{https://doi.org/10.1137/1115049}{doi:10.1137/1115049}.

\bibitem{Simon1993}
B.~Simon,
\emph{The Statistical Mechanics of Lattice Gases, Volume I},
Princeton University Press, 1993, Chapter V.

\bibitem{OS1}
K.~Osterwalder and R.~Schrader,
\emph{Axioms for Euclidean Green's functions},
Communications in Mathematical Physics \textbf{31} (1973), 83--112.
\href{https://doi.org/10.1007/BF01645738}{doi:10.1007/BF01645738}.

\bibitem{OS2}
K.~Osterwalder and R.~Schrader,
\emph{Axioms for Euclidean Green's functions II},
Communications in Mathematical Physics \textbf{42} (1975), 281--305.
\href{https://doi.org/10.1007/BF01608978}{doi:10.1007/BF01608978}.

\bibitem{OsterwalderSeiler}
K.~Osterwalder and E.~Seiler,
\emph{Gauge field theories on a lattice},
Annals of Physics \textbf{110} (1978), 440--471.
\href{https://doi.org/10.1016/0003-4916(78)90039-8}{doi:10.1016/0003-4916(78)90039-8}.

\bibitem{FILS}
J.~Fr\"ohlich, R.~Israel, E.~H.~Lieb, and B.~Simon,
\emph{Phase transitions and reflection positivity. I. General theory and
long range lattice models},
Communications in Mathematical Physics \textbf{62} (1978), 1--34.
\href{https://doi.org/10.1007/BF01940327}{doi:10.1007/BF01940327}.

\end{thebibliography}

\end{document}
